\begin{abstract}
Text-based emotion detection from social media is increasingly important because platforms such as Twitter generate large volumes of emotionally expressive informal text. This study presents MindTrace, a controlled comparative evaluation of four classification models --- SVM, XGBoost, CNN, and BiLSTM --- on a 416,809-sample Twitter corpus annotated across six emotion categories: anger, joy, sadness, fear, surprise, and love. Prior comparative studies vary in dataset selection, preprocessing, and evaluation, making direct performance comparisons unreliable. MindTrace addresses this gap by holding all experimental conditions constant across model families. A standardised NLP preprocessing pipeline performs tokenisation, negation-aware stopword removal, lemmatisation, and class balancing. ML models use TF-IDF vectorisation, whereas DL models use trainable embeddings. Performance is assessed using macro-F1 and confusion matrix analysis to account for substantial class imbalance (imbalance ratio $\approx 9.5{:}1$). BiLSTM achieved the highest test accuracy (94.70\%, macro-F1 0.9464), followed by CNN (93.13\%), SVM (91.80\%), and XGBoost (80.93\%). Confusion matrix analysis showed that Joy$\leftrightarrow$Love and Fear$\rightarrow$Surprise were the dominant misclassification pairs across all architectures, indicating that semantic overlap is the primary classification bottleneck. These findings provide practical guidance for architecture selection in emotion classification systems for social media analytics, mental health monitoring, and emotion-aware AI.
\end{abstract}

\noindent\textbf{Keywords:} emotion classification, natural language processing, machine learning, deep learning, BiLSTM, CNN, SVM, XGBoost, TF-IDF, Twitter, multi-class classification, social media analytics
